% TEMPLATE-MILC-v3.tex - plantilla maestra UNICA de las guias MILC v3.
%
% La usan las cuatro familias de guias, cada una con su driver:
%   · Grados 6-11          -> scripts/build-guias-g11.py
%   · Bebras               -> scripts/build-guias-bebras.py
%   · Semillero            -> scripts/build-guias-semillero.py
%   · Territorio interior  -> scripts/build-guias-territorio-interior.py
%
% Antes habia tres copias casi identicas (~400 lineas iguales, 6 distintas):
% cada mejora habia que aplicarla tres veces, y en la practica no se hacia
% (el motor de recursos vivia solo en dos de ellas). Lo que varia entre
% familias esta parametrizado en 6 placeholders que cada driver llena
% (se nombran aqui SIN su sintaxis real: el builder reemplaza por texto plano,
% tambien dentro de los comentarios, y un valor multilinea rompe el preambulo):
%
%   HEADER_IZQ        encabezado izquierdo de pagina
%   HEADER_DER        encabezado derecho de pagina
%   PORTADA_ETIQUETA  rotulo sobre el numero grande (GRADO/MOMENTO/MODULO)
%   PORTADA_SERIE     linea de serie de la portada
%   PORTADA_ID        identificador de la guia en la portada
%   PORTADA_CIRCULO   el circulito de grado (°), o vacio si no aplica
%   PDF_TITULO        titulo en texto plano (sin \textbf ni comillas TeX) para
%                     los metadatos del PDF (/Title); lo produce lib_xelatex.texto_pdf
%
% Composicion (auditoria 2026-09-03): las bandas de estacion piden espacio
% (\Needspace*) y prohiben el salto justo despues (\nopagebreak), asi no quedan
% huerfanas al pie; las cajas largas (softbox, drawbox, infoband, puentebox) son
% `breakable`, asi una caja de media pagina ya no arrastra todo a la siguiente
% dejando la anterior al 40 %. Todo texto sobre mostaza o verde va en milcNegro
% (blanco sobre #E5B400 da 1,93:1 y sobre #58A923 2,95:1; ninguno pasa WCAG AA).
% El resto de claves las documenta content/guias/_SCHEMA.md. El script valida
% que no queden placeholders sin reemplazar antes de escribir el archivo final.

% Etiquetado PDF (latex-lab): /Lang, estructura y texto alternativo de las
% figuras, para lectores de pantalla. Debe ir ANTES de \documentclass.
\DocumentMetadata{lang=es-CO,tagging=on}
\documentclass[11pt,letterpaper]{article}
% headheight=24pt: la cabecera derecha lleva dos lineas (identidad de la guia
% y estacion actual). headsep se acorta en la misma medida para que el bloque
% de texto quede exactamente donde estaba con headheight=12pt.
\usepackage[margin=2.0cm,headheight=24pt,headsep=13pt]{geometry}
\usepackage[table]{xcolor}
\usepackage{fontspec}
\usepackage[spanish]{babel}
\usepackage{tikz}
% Librerías TikZ para los diagramas de las guías (bloque `recursos:`):
%  · babel  → babel en español vuelve activos caracteres como «>», y TikZ los usa
%             en sus opciones (>=stealth, ->). Sin esta librería, cualquier
%             diagrama con flechas revienta con «\language@active@arg> has an
%             extra }». Debe ir ANTES de que se dibuje nada.
%  · positioning        → sintaxis «below left=2pt and 3pt».
%  · decorations.pathreplacing → llaves ({brace}) para señalar rangos.
\usetikzlibrary{babel,positioning,decorations.pathreplacing}
\usepackage{graphicx}
% Circuitos: el trabajo de electrónica se hace hoy con micro:bit y sus sensores
% integrados (MakeCode, sin cableado), así que no se necesitan esquemáticos.
% Si algún día entran montajes con protoboard, basta descomentar esta línea:
% los diagramas .tex/.tikz declarados en `recursos:` ya se incrustan solos.
% \usepackage{circuitikz}
\usepackage[most]{tcolorbox}
\usepackage{tabularx}
\usepackage{array}
\usepackage{fancyhdr}
\usepackage{titlesec}
\usepackage[document]{ragged2e}  % bandera a la derecha en todo el cuerpo
\usepackage{enumitem}
\usepackage{microtype}
\usepackage{etoolbox}
\usepackage{needspace}
% hyperref al final del preambulo: metadatos del PDF (titulo, autor, idioma,
% familia), marcadores por estacion (\pdfbookmark) y enlaces sin recuadro.
\usepackage[hidelinks,
  pdftitle={¿Qué es técnica? — del cuerpo al instrumento},
  pdfauthor={PhD. Álvaro Cárdenas Orozco},
  pdfsubject={Tecnología e Informática · Grado 9 · Guía 1 · MILC v3},
  pdflang={es-CO},
  bookmarksopen=true,bookmarksnumbered=false]{hyperref}

% Helvetica Neue no trae la flecha U+2192, asi que cada «→» escrito literalmente
% en el YAML se compilaba a NADA, en silencio (462 flechas en 61 guias: rutas del
% tipo «paleta → tipografia → lockup» salian rotas). Mapeamos el caracter a su
% version matematica, que si tiene glifo. Asi el autor puede seguir escribiendo
% «→» comodamente en el YAML.
\usepackage{newunicodechar}
\newunicodechar{→}{\ensuremath{\rightarrow}}
% Mismo problema con los simbolos de marca y vinetas: el «✓» de «una palomita ✓»
% salia como cuadrito vacio en el PDF de todas las guias que lo usaban.
\usepackage{pifont}
\newunicodechar{✓}{\ding{51}}
\newunicodechar{✗}{\ding{55}}
\newunicodechar{✔}{\ding{52}}
\newunicodechar{•}{\textbullet}
\newunicodechar{≥}{\ensuremath{\geq}}
\newunicodechar{≤}{\ensuremath{\leq}}

\setmainfont{Helvetica Neue}
\setsansfont{Helvetica Neue}
\setlength{\parindent}{0pt}
\setlength{\parskip}{6pt}
% Sin particion de palabras: en 6.º-7.º una palabra cortada por guion al final
% de linea («Comuni-cacion») frena la lectura mucho mas que una linea corta.
\hyphenpenalty=10000
\exhyphenpenalty=10000
\pretolerance=2000
\tolerance=3000
\setlength{\emergencystretch}{3em}
\linespread{1.10}
\renewcommand{\arraystretch}{1.25}
\setlist{leftmargin=5mm,itemsep=1mm,topsep=1mm}
\newcolumntype{Y}{>{\RaggedRight\arraybackslash}X}

\definecolor{milcMagenta}{HTML}{E6007E}
\definecolor{milcVino}{HTML}{5A0038}
\definecolor{milcTurquesa}{HTML}{0093A5}
\definecolor{milcVerde}{HTML}{58A923}
\definecolor{milcMostaza}{HTML}{E5B400}
\definecolor{milcOcre}{HTML}{B86600}
\definecolor{milcNegro}{HTML}{241816}
\definecolor{milcGris}{HTML}{F7F7F4}
\definecolor{milcLinea}{HTML}{E8E2DD}
\definecolor{milcGradePrimary}{HTML}{7C3AED}
\definecolor{milcGradeDark}{HTML}{4C1D95}
\definecolor{milcGradeSoft}{HTML}{EDE9FE}
\definecolor{milcGradeAccent}{HTML}{CFE7FF}
\definecolor{milcGradeText}{HTML}{FFFFFF}
\color{milcNegro}

\pagestyle{fancy}
\fancyhf{}
% Cabecera de dos lineas: identidad de la guia arriba y, a la derecha y en
% gris, la estacion en curso (\markright desde \milcestacion). Antes de la
% primera estacion la segunda linea va vacia; el \strut mantiene la altura
% para que la linea de arriba no baile entre paginas.
\lhead{\small Tecnología e Informática · Grado 9\\[-1pt]{\footnotesize\strut}}
\rhead{\small Guía 1 · MILC v3\\[-1pt]{\footnotesize\strut\color{milcNegro!65}\rightmark}}
\cfoot{\small\thepage}
\renewcommand{\headrulewidth}{0pt}

% Sobre mostaza (#E5B400) y verde (#58A923) el texto va en milcNegro; sobre el
% resto de la paleta, en blanco. Se usa dentro de fonttitle/fontupper, donde un
% \color posterior manda sobre coltitle.
\newcommand{\milctintasobre}[1]{%
  \ifboolexpr{test{\ifstrequal{#1}{milcMostaza}} or test{\ifstrequal{#1}{milcVerde}}}%
    {\color{milcNegro}}{\color{white}}}

\newtcolorbox{titlebox}[1]{
  enhanced,colback=#1,colframe=#1,arc=7mm,boxrule=0pt,
  left=7mm,right=7mm,top=3mm,bottom=3mm,
  before skip=8pt,after skip=8pt,
  fontupper=\sffamily\bfseries\Large\color{white}
}

\newtcolorbox{softbox}[3]{
  enhanced,breakable,pad at break=3mm,
  colback=#2,colframe=#1,borderline west={4pt}{0pt}{#1},
  arc=2mm,boxrule=.4pt,left=4mm,right=4mm,top=3mm,bottom=3mm,
  before skip=6pt,after skip=8pt,title={#3},coltitle=white,
  colbacktitle=#1,fonttitle=\sffamily\bfseries\milctintasobre{#1},fontupper=\RaggedRight,
  attach boxed title to top left={xshift=3mm,yshift=-2mm},
  boxed title style={arc=2mm, boxrule=0pt}
}

\newtcolorbox{drawbox}[2]{
  enhanced,breakable,pad at break=4mm,
  colback=white,colframe=#1,arc=4mm,boxrule=1pt,
  left=5mm,right=5mm,top=4mm,bottom=4mm,before skip=8pt,after skip=8pt,
  title={#2},coltitle=white,colbacktitle=#1,fonttitle=\sffamily\bfseries\milctintasobre{#1},
  fontupper=\RaggedRight,
  attach boxed title to top left={xshift=4mm,yshift=-2mm},
  boxed title style={arc=3mm, boxrule=0pt}
}

\newtcolorbox{infoband}[1]{
  enhanced,breakable,pad at break=2mm,colback=#1!8,colframe=#1!50,arc=2mm,boxrule=.5pt,
  left=4mm,right=4mm,top=2mm,bottom=2mm,before skip=4pt,after skip=4pt,
  fontupper=\small\RaggedRight
}

% Puente narrativo entre secciones (contrato editorial Conéctate).
% Da continuidad: "Acabas de X. Ahora vas a Y porque Z."
% Visualmente: banda izquierda en color del grado, texto en itálica pequeña.
\newtcolorbox{puentebox}{
  enhanced,breakable,pad at break=2mm,colback=milcGradeSoft,colframe=milcGradeDark,
  borderline west={3pt}{0pt}{milcGradeDark},
  arc=0mm,boxrule=0pt,left=5mm,right=4mm,top=2.5mm,bottom=2.5mm,
  before skip=4pt,after skip=8pt,
  fontupper=\itshape\small\color{milcNegro}\RaggedRight
}
% La flecha va en modo matematico a proposito: Helvetica Neue NO tiene el glifo
% U+2192, asi que \textrightarrow se compilaba a NADA («Missing character: There
% is no → in font Helvetica Neue») y el puente salia sin flecha. En modo
% matematico la flecha viene de la fuente matematica y si se dibuja.
\newcommand{\puente}[1]{%
  \begin{puentebox}%
    \textcolor{milcGradeDark}{$\rightarrow$}\hspace{2mm}#1%
  \end{puentebox}%
}

\newcommand{\linead}[1]{\noindent\color{milcLinea}\rule{#1}{0.5pt}\color{milcNegro}}
\newcommand{\respuesta}[1]{\par\vspace{2mm}\foreach \n in {1,...,#1}{\noindent\color{milcLinea}\rule{\linewidth}{0.5pt}\par\vspace{5mm}}\color{milcNegro}}
\newcommand{\checkbox}{\textcolor{milcGradeDark}{\fbox{\phantom{x}}}\hspace{5pt}}

% ─── Listas aireadas para las guías socioemocionales ────────────────────
% Componer las verificaciones como «\checkbox texto\\» deja el texto que salta
% de línea debajo del cuadrito y apelmaza el bloque. Estos entornos alinean la
% continuación y airean los ítems. Los usa Territorio Interior; cualquier
% familia puede hacerlo.
\newlist{checklista}{itemize}{1}
\setlist[checklista]{
  label=\textcolor{milcGradeDark}{\fbox{\phantom{x}}},
  leftmargin=9mm, labelsep=3.2mm, itemsep=2.4mm,
  topsep=2.5mm, parsep=0pt, partopsep=0pt, before=\RaggedRight
}
\newlist{pasoslista}{enumerate}{1}
\setlist[pasoslista]{
  label=\textcolor{milcGradeDark}{\sffamily\bfseries\arabic*.},
  leftmargin=8mm, labelsep=3mm, itemsep=2.8mm,
  topsep=1mm, parsep=0pt, partopsep=0pt, before=\RaggedRight
}

\newcommand{\phasecard}[4]{
\begin{tcolorbox}[enhanced,colback=#3,colframe=#2,arc=3mm,boxrule=.5pt,height=3.05cm,valign=center,halign=center,left=1.5mm,right=1.5mm,top=2mm,bottom=2mm]
{\sffamily\bfseries\fontsize{22}{24}\selectfont\textcolor{#2}{#1}}\par
{\sffamily\bfseries\small #4}
\end{tcolorbox}
}

% ── Piezas del rediseno 2026 (legibilidad 6.º-11.º) ───────────────────
% Banda de estacion: reemplaza el titlebox macizo. Titulo a la izquierda,
% dato practico (tiempo/modalidad) a la derecha, en una sola linea.
% Nunca huerfana: pide 9 lineas libres (o salta de pagina rellenando con
% \vfill) y prohibe el salto justo despues de la banda. Nueve y no seis:
% la caja partible que sigue necesita ~80pt (titulo adosado, relleno y dos
% lineas) para arrancar en la misma pagina; con menos, tcolorbox la manda
% entera a la siguiente y la banda se queda sola. El penalty 10000 va
% en `after`, ANTES del espacio posterior: un `after skip` (aunque sea 0pt)
% es un \vskip tras la caja, y ese glue es un punto de corte legal que un
% \nopagebreak escrito despues de \end{tcolorbox} ya no cierra. Cada banda
% deja marcador PDF y actualiza la cabecera derecha.
\newcounter{milcestacion}
\newcommand{\milcestacion}[2]{%
\Needspace*{9\baselineskip}%
\stepcounter{milcestacion}%
\pdfbookmark[1]{#2}{milcestacion\themilcestacion}%
\markright{#2}%
\begin{tcolorbox}[enhanced,colback=#1,colframe=#1,arc=2mm,boxrule=0pt,
  left=5mm,right=5mm,top=2.6mm,bottom=2.6mm,before skip=11pt,
  after={\par\nopagebreak\vskip7pt}]
{\sffamily\bfseries\fontsize{15}{18}\selectfont\milctintasobre{#1}#2}
\end{tcolorbox}}

% Tarjeta de paso numerada: sustituye las tablas «Pilar / Aplicacion», que
% eran formato de consulta docente, no de estudiante. El numero va en un
% circulo sobre el borde para que la secuencia se lea de un vistazo.
\newcommand{\milcpaso}[3]{%
\Needspace{4\baselineskip}%
\begin{tcolorbox}[enhanced,colback=white,colframe=milcLinea,arc=1.5mm,boxrule=.9pt,
  left=14mm,right=4mm,top=3mm,bottom=3mm,before skip=4pt,after skip=4pt,
  fontupper=\RaggedRight,
  overlay={\node[anchor=north west,circle,fill=#2,text=white,inner sep=0pt,
    minimum size=7.4mm,font=\sffamily\bfseries\small]
    at ([xshift=3.6mm,yshift=-3.2mm]frame.north west) {#1};}]
#3
\end{tcolorbox}}

% Casilla marcable de tamano util con lapiz (la anterior era \fbox{\phantom{x}},
% de alto variable segun el texto que la seguia).
\newcommand{\milcmarca}{\framebox[3.8mm]{\rule{0pt}{3.4mm}}}

% Fila marcable de lista suelta (checks de la estacion 1).
\newcommand{\milccheckrow}[1]{%
\par\vspace{1.8mm}\noindent
\makebox[6.4mm][l]{\raisebox{-.55mm}{\textcolor{milcNegro}{\milcmarca}}}%
\parbox[t]{\dimexpr\linewidth-6.4mm\relax}{\RaggedRight #1}\par}

% Celda marcable dentro de la rubrica (tabla de autoevaluacion). Misma
% sangria francesa, pero pensada para vivir dentro de una columna X.
\newcommand{\milcopcion}[1]{%
\makebox[5.2mm][l]{\raisebox{-.55mm}{\textcolor{milcNegro}{\milcmarca}}}%
\parbox[t]{\dimexpr\linewidth-5.2mm\relax}{\RaggedRight #1}}

% ── Recursos visuales de la guía ──────────────────────────────────────
% Declarados en el bloque `recursos:` del YAML y emitidos por el builder.
% \guiaFigura   → imagen rasterizada o PDF (foto, carta celeste, montaje).
% \guiaDiagrama → fuente TikZ/circuitikz (.tex) incrustada como vector:
%                 es la vía para circuitos y esquemáticos editables.
% [#1] = texto alternativo (alt) · #2 = ruta relativa al .tex · #3 = pie
% El alt llega al PDF etiquetado (/Alt) y lo lee el lector de pantalla.
\newcommand{\guiaFigura}[3][]{%
\begin{center}
\includegraphics[width=0.88\linewidth,height=0.38\textheight,keepaspectratio,alt={#1}]{#2}\\[1.5mm]
{\footnotesize\itshape\textcolor{milcNegro!75}{#3}}
\end{center}
}

\newcommand{\guiaDiagrama}[2]{%
\begin{center}
\input{#1}\\[1.5mm]
{\footnotesize\itshape\textcolor{milcNegro!75}{#2}}
\end{center}
}

\begin{document}
\thispagestyle{empty}

% ── Portada ───────────────────────────────────────────────────────────
% Antes: un rectangulo de color a pagina completa. Se veia bien en pantalla,
% pero en la fotocopiadora del colegio se comia el toner y salia gris sucio.
% Ahora la banda ocupa el tercio superior y el resto va en blanco.
\begin{tikzpicture}[remember picture,overlay]
  \path[fill=milcGradePrimary]
    (current page.north west) rectangle ([yshift=-10.6cm]current page.north east);

  \node[anchor=north west,text=milcGradeText] at ([xshift=2.0cm,yshift=-1.55cm]current page.north west)
    {\sffamily\bfseries\fontsize{12}{14}\selectfont GRADO};
  \node[anchor=north west,text=milcGradeText,opacity=.85] at ([xshift=2.0cm,yshift=-2.35cm]current page.north west)
    {\sffamily\bfseries\fontsize{8.5}{10}\selectfont SERIE GUÍAS · TECNOLOGÍA E INFORMÁTICA};
  \node[anchor=north east,text=milcGradeText,opacity=.85,align=right] at ([xshift=-2.0cm,yshift=-1.55cm]current page.north east)
    {\sffamily\bfseries\fontsize{9}{12}\selectfont I.E. Sor María Juliana\\Cartago, Valle del Cauca};

  \node[anchor=west,text=milcGradeText] (gradonum) at ([xshift=1.85cm,yshift=-7.1cm]current page.north west)
    {\sffamily\bfseries\fontsize{112}{112}\selectfont 9};
  % El circulito de grado (°) solo aplica a las guias de grado («11°»); en
  % Bebras y Semillero el numero grande es un momento/modulo, no un grado.
  % Se posiciona relativo al nodo `gradonum`, no en coordenadas absolutas.
  \draw[milcGradeText,line width=1.6pt]
    ([xshift=2.0mm,yshift=-4.5mm]gradonum.north east) circle (.15cm);

  \node[anchor=west,text=milcGradeText,text width=11.4cm,align=left] at ([xshift=1.5cm,yshift=0.5cm]gradonum.east)
    {
     {\sffamily\bfseries\fontsize{9}{11}\selectfont\MakeUppercase{Período 1 · Historia de la técnica}}\\[3mm]
     {\sffamily\bfseries\fontsize{21}{25}\selectfont\setlength{\baselineskip}{27pt}¿Qué es técnica? ---\\[4pt]del cuerpo\\[4pt]al instrumento}\\[3.5mm]
     {\sffamily\bfseries\fontsize{15}{18}\selectfont Guía 1-9-TIC}
    };
\end{tikzpicture}

\vspace*{9.4cm}
\vfill

{\sffamily\bfseries\fontsize{8}{10}\selectfont\textcolor{milcGradeDark}{LO QUE TE LLEVAS HOY}}

\begin{tcolorbox}[enhanced,colback=white,colframe=milcNegro,arc=2mm,boxrule=1.6pt,
  left=5mm,right=5mm,top=4mm,bottom=4mm,before skip=3pt,after skip=12pt,
  fontupper=\RaggedRight\fontsize{11}{15.5}\selectfont]
Tu definición personal de técnica en una frase, acompañada de cinco ejemplos concretos de tu vida. Cada ejemplo dice qué función del cuerpo extiende el objeto y en qué etapa está: cuerpo, herramienta o máquina. Al menos uno viene de un oficio o de una tradición que conozcas de cerca. Cierra con la ficha de las tres etapas explicadas con tus palabras.
\end{tcolorbox}

\begin{tcolorbox}[enhanced,colback=milcGris,colframe=milcGris,arc=2mm,boxrule=0pt,
  left=5mm,right=5mm,top=3.5mm,bottom=3.5mm,after skip=12pt]
{\sffamily\bfseries\fontsize{8}{10}\selectfont\textcolor{milcNegro!55}{ESTA GUÍA ES DE}}\\[3mm]
\begin{tabularx}{\linewidth}{X p{3.2cm} p{3.2cm}}
\linead{\linewidth} & \linead{3.0cm} & \linead{3.0cm}\\[-1mm]
{\fontsize{8.5}{10}\selectfont\textcolor{milcNegro!55}{Nombre completo}} &
{\fontsize{8.5}{10}\selectfont\textcolor{milcNegro!55}{Grupo}} &
{\fontsize{8.5}{10}\selectfont\textcolor{milcNegro!55}{Fecha}}\\
\end{tabularx}
\end{tcolorbox}

\vfill

\noindent\textcolor{milcLinea}{\rule{\linewidth}{1pt}}\\[2mm]
{\sffamily\bfseries\fontsize{9.5}{12}\selectfont Autor: Dr. Álvaro Cárdenas Orozco}\\[1mm]
{\sffamily\fontsize{8.5}{11}\selectfont\textcolor{milcNegro!55}{Metodología MILC · apertura ancestral · pensamiento computacional · triángulo Dussel–estoicismo–Floridi}}

\newpage

\section*{Apertura: ¿Qué es técnica? --- del cuerpo al instrumento}

\begin{softbox}{milcGradeDark}{milcGradeSoft}{Saber ancestral}
Entre los embera, el jaibaná no nace jaibaná ni hereda el cargo: aprende de un maestro, al que acompaña durante años. El antropólogo Luis Guillermo Vasco Uribe lo documentó en \emph{Jaibanás: los verdaderos hombres} (1985), tras trabajar con los embera chamí de Risaralda y con los embera del río Garrapatas, en el Valle del Cauca. Se aprende estando al lado, mirando y haciendo, no recibiendo instrucciones. Por eso lo primero que se aprende de alguien es cómo mira el mundo. Fíjate en lo que eso dice de la técnica. Antes que cualquier objeto está un cuerpo que hace algo, y alguien que aprende mirándolo. La herramienta llega después. La \textbf{cara de exclusión}: el saber del jaibaná es restringido y ha sido objeto de extracción y exotización. Esta guía toma solo el modo de aprender, no el contenido, y ahí se detiene. Hoy vas a mirar tu propia vida con esa lente y a decidir qué llamas técnica.\par\smallskip{\footnotesize\textcolor{milcNegro!70}{\textbf{Fuente.} Vasco Uribe, L. G. (1985). \emph{Jaibanás: los verdaderos hombres}. Fondo de Promoción de la Cultura, Banco Popular.}}
\end{softbox}

\begin{softbox}{milcTurquesa}{milcGris}{Saber contemporáneo}
Toda técnica humana pasa por tres etapas que se acumulan, no se reemplazan. \textbf{Cuerpo}: el órgano hace el trabajo. La mano corta, la voz cuenta, la memoria guarda. \textbf{Herramienta}: un objeto extiende esa función. El cuchillo, el ábaco, el cuaderno. \textbf{Máquina}: un sistema automatiza el trabajo con una fuente de energía. El trapiche, la calculadora, el computador. Ninguna etapa borra a la anterior: sigues cortando con la mano aunque exista la licuadora. Y hay una distinción que sí importa. Una técnica \textbf{amplía} el cuerpo cuando te deja hacer más de lo que hacías. Lo \textbf{reemplaza} cuando deja de hacer falta que tú sepas. El GPS amplía tu orientación si lo lees; la reemplaza si solo lo obedeces. Ninguna de las dos es mala por sí sola. Lo que no conviene es no saber en cuál estás.
\end{softbox}

\begin{softbox}{milcMostaza}{milcGris}{Pregunta-puente}
\textit{El aprendiz de jaibaná pasa años mirando a alguien hacer. ¿Qué sabes hacer tú porque viste a alguien hacerlo, y qué sabes hacer solo porque un aparato lo hace por ti?}
\end{softbox}

\begin{softbox}{milcVerde}{milcGris}{Saber y saber hacer}
Al terminar podrás: (1) \textbf{identificar} en tu entorno cinco objetos que extienden una función concreta de tu cuerpo; (2) \textbf{explicar} con tus palabras la diferencia entre cuerpo, herramienta y máquina, con un ejemplo local de cada una; (3) \textbf{crear} tu definición personal de técnica en una frase, probada con cinco ejemplos de tu vida.
\end{softbox}



\small
\begin{tabularx}{\textwidth}{>{\bfseries}p{3.0cm}Y}
\rowcolor{milcGradePrimary!12}Autor & Dr. Álvaro Cárdenas Orozco\\
DBA & Análisis crítico de la historia de la tecnología y su impacto social (MEN, T\&I)\\
Referentes & Dussel (técnica situada) · Estoicismo (téchne del cuerpo) · Floridi (infoesfera)\\
Producto final & Tu definición personal de técnica en una frase, acompañada de cinco ejemplos concretos de tu vida. Cada ejemplo dice qué función del cuerpo extiende el objeto y en qué etapa está: cuerpo, herramienta o máquina. Al menos uno viene de un oficio o de una tradición que conozcas de cerca. Cierra con la ficha de las tres etapas explicadas con tus palabras.\\
\end{tabularx}
\normalsize

\puente{Este año recorres la historia de la técnica, del primer filo de piedra al chip. Hoy pones el punto de partida: qué entiendes tú por técnica, antes de que nadie te lo defina. Vas a comparar esa frase con la que tengas en diciembre.}

\milcestacion{milcGradeDark}{Ruta MILC: así vamos a aprender}
\begin{tabularx}{\textwidth}{>{\centering\arraybackslash}X>{\centering\arraybackslash}X>{\centering\arraybackslash}X>{\centering\arraybackslash}X}
\phasecard{1}{milcVerde}{milcGris}{Escucha\\[-1mm]\small Observo, escucho y pregunto.} &
\phasecard{2}{milcTurquesa}{milcGris}{Entender\\[-1mm]\small Sistematizo y ordeno.} &
\phasecard{3}{milcMagenta}{milcGris}{Hacer\\[-1mm]\small Construyo y aplico.} &
\phasecard{4}{milcVino}{milcGris}{Cierre\\[-1mm]\small Evalúo y reflexiono.}
\end{tabularx}


\puente{Antes de la teoría vas a mirar lo que tienes alrededor. Cinco objetos, y qué hace cada uno que tu cuerpo solo no puede. Sin esa observación, cualquier idea de técnica que tengas es prestada.}

\milcestacion{milcVerde}{Estación 1 de 3 · Escucha}

\begin{softbox}{milcVerde}{milcGris}{Escena para escuchar}
\textbf{Actividad 1 · IDENTIFICA --- Técnicas en tu cotidianidad} (15 min · individual). (1) Sin moverte de donde estás, encuentra cinco objetos que extiendan tu cuerpo. Pueden ser tan obvios como un lápiz o tan invisibles como la silla. (2) Escribe para cada uno qué función corporal concreta amplía. No vale «ayudar»: di si es fuerza, alcance, memoria, vista u oído. (3) Marca cuáles ya no notas y cuáles todavía sientes nuevos. (4) Señala al menos uno que venga de un oficio o de una tradición, no de una fábrica.
\end{softbox}

{\sffamily\bfseries\fontsize{8}{10}\selectfont\textcolor{milcNegro!55}{MARCA CUANDO LO HAYAS HECHO}}
\milccheckrow{Tengo cinco objetos con su función corporal concreta.}
\milccheckrow{Marqué cuáles ya no noto y cuáles todavía siento nuevos.}
\milccheckrow{Señalé al menos uno que viene de un oficio o una tradición.}

\begin{infoband}{milcVerde}


\textbf{Lo que va al cuaderno (Actividad 1 · IDENTIFICA).} \textbf{Título:} «Actividad 1 --- Técnicas en mi cotidianidad». \textbf{Formato:} tabla de 5 filas y 3 columnas (objeto / qué del cuerpo extiende / lo noto o ya es cuerpo). \textbf{Extensión:} un tercio de página. \textbf{Sabes que terminaste cuando} ninguna fila dice «ayudar» y una está marcada como de oficio o tradición.
\end{infoband}

\puente{Ya observaste. Ahora vas a ordenar lo observado en tres etapas y a probar que las entendiste dando un ejemplo local de cada una.}

\milcestacion{milcTurquesa}{Estación 2 de 3 · Entender}

\begin{softbox}{milcTurquesa}{milcGris}{Lo que vas a entender}
\textbf{Actividad 2 · EXPLICA --- Cuerpo, herramienta, máquina} (20 min · parejas). (1) Con tu pareja, escriban las tres etapas con una frase propia cada una. (2) Pongan un ejemplo de su vida diaria en cada etapa. (3) Pongan un ejemplo de un oficio o una tradición que conozcan en cada etapa. (4) Tomen uno de sus objetos de la Actividad 1 y decidan si amplía o reemplaza, con una razón escrita.
\end{softbox}

\milcpaso{1}{milcTurquesa}{\textbf{Las tres etapas se acumulan.} El cuerpo hace. La herramienta extiende. La máquina automatiza con energía propia. Ninguna borra a la anterior: sigues contando con los dedos aunque tengas calculadora.}
\milcpaso{2}{milcTurquesa}{\textbf{Cómo se distingue una máquina de una herramienta.} La herramienta necesita tu fuerza todo el tiempo. La máquina necesita otra fuente de energía y repite el gesto sin ti. Un trapiche es máquina aunque no tenga un solo cable.}
\milcpaso{3}{milcTurquesa}{\textbf{Ampliar o reemplazar.} Amplía cuando te deja hacer más de lo que hacías. Reemplaza cuando deja de hacer falta que tú sepas. No es que una sea buena y otra mala: es que conviene saber en cuál estás parado.}
\milcpaso{4}{milcTurquesa}{\textbf{Cuatro preguntas para cualquier objeto.} Qué función del cuerpo extiende. En qué etapa está. Si amplía o reemplaza. Qué cambiaría en tu día si mañana no existiera.}

\begin{drawbox}{milcTurquesa}{Cuerpo, herramienta, máquina}
\textbf{Cuerpo:} el órgano hace el trabajo. \\[1mm] \textbf{Herramienta:} extiende una función y necesita tu fuerza. \\[1mm] \textbf{Máquina:} automatiza con otra fuente de energía. \\[1mm] \textbf{Prueba:} ¿amplía lo que sé hacer o hace que no me haga falta saberlo?
\end{drawbox}

\begin{softbox}{milcMostaza}{milcGris}{Errores comunes y su corrección}
(1) \textbf{Creer que técnica es solo lo electrónico}: una silla también es técnica. (2) \textbf{Confundir máquina con aparato digital}: un molino de agua es máquina y no tiene un solo circuito. (3) \textbf{Ordenar las etapas como si fueran mejores unas que otras}: son distintas, no superiores. (4) \textbf{Describir la función como «ayudar»}: eso no dice nada; di fuerza, alcance, memoria, vista u oído. (5) \textbf{Olvidar que toda técnica resuelve algo}: si no resuelve un problema, es adorno.
\end{softbox}

\begin{infoband}{milcTurquesa}


\textbf{Lo que va al cuaderno (Actividad 2 · EXPLICA).} \textbf{Título:} «Actividad 2 --- Cuerpo, herramienta, máquina». \textbf{Formato:} tres fichas, una por etapa, cada una con la definición propia, un ejemplo de tu vida y un ejemplo de un oficio o tradición. \textbf{Extensión:} media página. \textbf{Sabes que terminaste cuando} puedes explicar las tres etapas sin mirar las fichas.
\end{infoband}

\puente{Tienes el mapa. Toca tomar posición y escribir tu propia frase, con cinco ejemplos que la sostengan.}

\milcestacion{milcMagenta}{Estación 3 de 3 · Hacer}

\begin{softbox}{milcMagenta}{milcGris}{Cómo vas a construir tu producto}
\textbf{Actividad 3 · CREA --- Tu definición personal de técnica} (30 min · individual). (1) Responde primero en voz baja: si tuvieras que explicarle a un niño qué es la técnica, ¿qué le dirías? (2) Escribe esa respuesta en una sola frase. (3) Lista cinco objetos de tu vida que la prueben, y di en cada uno qué función del cuerpo extiende y en qué etapa está. (4) Asegúrate de que al menos uno venga de un oficio o una tradición. (5) Quita de tu frase toda palabra que no haga falta. (6) Léesela a un compañero y pídele que te diga con sus palabras qué entendió; anota si coincide.
\end{softbox}

\milcpaso{1}{milcMagenta}{\textbf{Tu entrega tiene dos piezas.} Una frase con tu definición. Cinco ejemplos numerados, cada uno con su función corporal y su etapa.}
\milcpaso{2}{milcMagenta}{\textbf{Una frase, no un párrafo.} Si tu definición necesita tres renglones, todavía no la entendiste. La prueba es leerla en voz alta de un tirón, sin corregir a mitad de camino.}
\milcpaso{3}{milcMagenta}{\textbf{Cada ejemplo dice una sola cosa.} Si el celular te sirve para recordar, comunicar y distraerte, elige una función y escríbela. Cinco ejemplos claros valen más que diez amontonados.}
\milcpaso{4}{milcMagenta}{\textbf{Al menos un ejemplo de oficio o tradición.} No por folclor, sino porque obliga a mirar técnicas que no vienen en un empaque. El molde de una modista o el mango gastado de una herramienta enseñan lo mismo que un aparato.}

\begin{drawbox}{milcMagenta}{Antes de entregar}
\checkbox La definición cabe en una frase y se lee de un tirón. \\[1mm] \checkbox Está escrita con tus palabras, no copiada. \\[1mm] \checkbox Cinco ejemplos concretos de tu vida, no genéricos. \\[1mm] \checkbox Cada ejemplo dice qué función del cuerpo extiende y en qué etapa está. \\[1mm] \checkbox Al menos uno viene de un oficio o una tradición.
\end{drawbox}

\begin{softbox}{milcMagenta}{milcGris}{Plantilla de trabajo}
\textbf{Mi definición.} \textit{Una frase, sin palabras de más.} \\[1mm] \textbf{Ejemplo 1.} \textit{Objeto, función corporal, etapa.} \\[1mm] \textbf{Ejemplos 2 a 5.} \textit{Igual, variando de etapa y mezclando lo cotidiano con lo de oficio.} \\[1mm] \textbf{Prueba.} \textit{Leerla en voz alta a alguien más pequeño.}
\end{softbox}

\begin{infoband}{milcMagenta}


\textbf{Lo que va al cuaderno (Actividad 3 · CREA).} \textbf{Título:} «Actividad 3 --- Mi definición personal de técnica». \textbf{Formato:} la frase con tu definición y los cinco ejemplos numerados, cada uno con su función corporal y su etapa. \textbf{Extensión:} media página. \textbf{Sabes que terminaste cuando} alguien más pequeño que tú entendió la frase a la primera.
\end{infoband}

\puente{Lo que entregas hoy: una frase tuya y cinco ejemplos. La prueba de calidad: un primo más pequeño entiende tu definición sin que se la expliques dos veces.}

\milcestacion{milcMostaza}{Producto de la sesión}
\textbf{Definición personal de técnica en una frase + cinco ejemplos con función y etapa}.
\begin{softbox}{milcMostaza}{milcGris}{Criterios de calidad}
(1) La definición cabe en una frase y está en palabras tuyas. (2) Un compañero la leyó y dijo qué entendió sin que se la explicaras. (3) Los cinco ejemplos son objetos concretos de tu vida. (4) Cada ejemplo nombra la función corporal, no dice solo «ayudar». (5) Cada ejemplo dice en qué etapa está, cuerpo, herramienta o máquina. (6) Al menos un ejemplo viene de un oficio o una tradición.
\end{softbox}

\puente{Antes de cerrar, mira lo que hiciste desde cinco lados. Definir qué es la técnica decide cómo vas a mirar todo lo demás del año.}

\milcestacion{milcVino}{Evaluación liberadora · cinco dimensiones}

\begin{softbox}{milcVino}{milcGris}{Sentido de la evaluación}
Evaluar no es solo revisar una nota. Es reconocer qué aprendí, qué cambió en mí, qué emociones manejé, cómo me relaciono con la ciudad y la comunidad, y qué le dejaré a quien viene detrás.
\end{softbox}

\textbf{Mi reflexión liberadora en cinco dimensiones.} Completa cada frase iniciada:

\small
\begin{tabularx}{\textwidth}{>{\bfseries\RaggedRight\arraybackslash}p{3.4cm}Y>{\RaggedRight\arraybackslash}p{4.6cm}}
\rowcolor{milcVino}\color{white}Dimensión & \color{white}\bfseries Mi respuesta (frame starter) & \color{white}\bfseries Algo que descubrí\\
1. Personal & Lo que más cambió en mí fue \linead{6.5cm} & \linead{4.2cm}\\
2. Emocional & La emoción más fuerte fue \linead{2.5cm} y la regulé \linead{3cm} & \linead{4.2cm}\\
3. Ciudadana & En democracia esto importa porque \linead{6cm} & \linead{4.2cm}\\
4. Local (Cartago) & En mi barrio sería útil para \linead{6.5cm} & \linead{4.2cm}\\
5. Intergeneracional & A mi abuela/abuelo le diría \linead{6.5cm} & \linead{4.2cm}\\
\end{tabularx}
\normalsize

\begin{infoband}{milcVino}
\textbf{Para la próxima sustentación.} Vuelve a esta tabla en seis meses. Verás que las respuestas cambian — no porque lo aprendido cambie, sino porque tú cambias. La evaluación liberadora es una conversación contigo a través del tiempo.
\end{infoband}


\puente{Para terminar, tres citas verificadas sobre hacer y saber: Dussel sobre la técnica que no es teoría aplicada, Séneca sobre enseñar a hacer y no a decir, Floridi sobre el tiempo que las tecnologías ahorran y luego matan.}

\milcestacion{milcGradeDark}{Triángulo de pensamiento · cierre filosófico}

\begin{softbox}{milcVino}{milcGris}{Dussel · lente del nosotros}
\textit{«La tecnología no es ciencia aplicada (teoría concretada), sino que es técnica científica.»}

{\footnotesize\itshape\color{milcNegro!70}--- Enrique Dussel · Filosofía de la liberación (1977), §4.3.2.4}

La técnica no salió del laboratorio para bajar a la mano. Salió de la mano y después se dejó explicar. El aprendiz de jaibaná no recibe primero la teoría; está al lado, mira y hace. Lo mismo pasa con casi todo lo que sabes hacer bien.

\textbf{¿Qué sé hacer sin poder explicarlo del todo, y quién me lo enseñó estando al lado?}
\end{softbox}
\linead{\linewidth}\\[1mm]
\linead{\linewidth}

\begin{softbox}{milcMostaza}{milcGris}{Estoicismo · Séneca · Cartas a Lucilio, 20 (c. 64 d.C.) · cuidado interior}
\textit{«La filosofía enseña a hacer, no a decir; y exige que cada uno viva según su propia ley, que la vida no desdiga de la palabra. (trad. propia)»}

{\footnotesize\itshape\color{milcNegro!70}--- Séneca · Cartas a Lucilio, 20 (c. 64 d.C.)}

Tu definición de hoy se prueba en los cinco ejemplos, no en lo bonito que suene. Si la frase dice una cosa y tus ejemplos muestran otra, la frase está de adorno. Escribir lo que uno de verdad hace es más difícil que escribir lo que suena bien.

\textbf{¿Mi definición describe cómo uso las cosas de verdad, o cómo me gustaría usarlas?}
\end{softbox}
\linead{\linewidth}\\[1mm]
\linead{\linewidth}

\begin{softbox}{milcTurquesa}{milcGris}{Floridi · lente de la infoesfera}
\textit{«Las tecnologías se usan primero para ahorrar tiempo y luego para matarlo. (trad. propia)»}

{\footnotesize\itshape\color{milcNegro!70}--- Luciano Floridi · Commentary on the Onlife Manifesto (2015), § 3.1}

Casi todo objeto de tu lista entró a tu vida prometiendo ahorrarte algo. Vale la pena revisar qué hiciste con lo que te ahorró. A veces el tiempo ganado se usó en otra cosa; a veces se lo quedó el mismo aparato que lo liberó.

\textbf{¿En qué gasté el tiempo que me ahorró el objeto que más uso?}
\end{softbox}
\linead{\linewidth}\\[1mm]
\linead{\linewidth}

\begin{infoband}{milcNegro}
\textbf{Nota para el docente.} Las tres son citas textuales y llevan su referencia debajo. Puedes remitir a la obra en clase.
\end{infoband}

\milcestacion{milcVino}{Mi autoevaluación · marca una casilla por fila}

{\small
\setlength{\extrarowheight}{2.2pt}
\begin{tabularx}{\textwidth}{>{\RaggedRight\arraybackslash\bfseries}p{3.0cm}YYY}
\rowcolor{milcVino}\color{white}\bfseries Criterio & \color{white}\bfseries Lo logré & \color{white}\bfseries En proceso & \color{white}\bfseries Necesito apoyo\\[.6mm]
Comprensión del tema & \milcopcion{Explico las ideas con mis palabras.} & \milcopcion{Reconozco las ideas, falta precisión.} & \milcopcion{Necesito repasar lo básico.}\\[.8mm]
\rowcolor{milcGris}Pensamiento computacional & \milcopcion{Usé los pilares al armar mi producto.} & \milcopcion{Usé algunos pilares.} & \milcopcion{Necesito releer la estación 2.}\\[.8mm]
Aplicación y producto & \milcopcion{Mi producto cumple los criterios.} & \milcopcion{Está armado, con ajustes pendientes.} & \milcopcion{Aún está en construcción.}\\[.8mm]
\rowcolor{milcGris}Reflexión 5 dimensiones & \milcopcion{Respondí cada una con honestidad.} & \milcopcion{Respondí algunas con detalle.} & \milcopcion{Mis respuestas son muy generales.}\\[.8mm]
Triángulo de pensamiento & \milcopcion{Conecté las 3 voces con mi trabajo.} & \milcopcion{Conecté al menos una.} & \milcopcion{No las relaciono con mi proceso.}\\[.8mm]
\end{tabularx}}

\vspace{3mm}
\textbf{Hoy entendí que:}\\[1mm]
\linead{\linewidth}\\[1mm]
\linead{\linewidth}

\textbf{Mi compromiso para la próxima sesión es:}\\[1mm]
\linead{\linewidth}\\[1mm]
\linead{\linewidth}

\begin{softbox}{milcMostaza}{milcGris}{Cita de despedida · Marco Aurelio}
\textit{``El obstáculo en el camino se vuelve el camino. Lo que estorba al camino se vuelve camino.''}\\[1mm]
Cada error de hoy, bien observado, se convierte en pieza del camino que pisarás mañana.
\end{softbox}

\small
\begin{tabularx}{\textwidth}{>{\bfseries}p{3.6cm}Y}
\rowcolor{milcGradeDark!12}Nombre completo: & \linead{10cm}\\
Fecha: & \linead{4cm} \quad Curso/grupo: \linead{4cm}\\
Mi firma: & \linead{9cm}\\
\end{tabularx}
\normalsize

\begin{infoband}{milcGradeDark}
\textbf{Para la próxima sesión.} Lleva el cuaderno contigo tres días. Cada vez que uses un objeto técnico, anota una línea: nombre y qué del cuerpo extiende. Al tercer día lee la lista completa. En la sesión 2 vas a estudiar máquinas simples del campo colombiano, y varias van a estar ahí.
\end{infoband}

\end{document}
